\chapter{Kết luận và hướng phát triển}
\label{chap:ket-luan}

\section{Kết luận}
\label{sec:ket-luan}

Tiểu luận đặt câu hỏi ``grounded hay guessed'' cho SemLoc~\cite{yang2026semloc}, phương pháp định vị lỗi ngữ nghĩa neo
lập luận của LLM vào ràng buộc thực thi được. Bằng cách tái hiện từ mã nguồn và dữ liệu của tác giả, kiểm chứng lại bộ
dữ liệu bằng thực thi, lặp lại thực nghiệm ba lần với kiểm định thống kê và phân rã đóng góp của từng tín hiệu, nhóm rút
ra ba kết luận.

Thứ nhất, kết quả công bố tái hiện được, nhưng ở hạng 1 phần lớn độ chính xác của SemLoc đến từ dòng lỗi mà LLM tự đoán
và được cộng điểm thưởng --- một thành phần không có trong bài báo: dòng đoán đóng góp
\DecSemlocRepairedAllAccOnePhiGuess{} trên \DecSemlocRepairedAllAccOneFull{} điểm Acc@1, và trên lỗi thật thì mọi lần
đúng ở hạng 1 đều là dòng đoán. Bằng chứng thực thi có ích thật, nhưng chủ yếu ở top-3.

Thứ hai, bằng chứng thực thi phát huy tốt nhất khi được đưa thẳng cho LLM: một prompt duy nhất có số kiểm thử đạt/không
đạt chạy qua mỗi dòng vượt cả pipeline nhiều bước ở cả Acc@1 và Acc@3, tốn ít token hơn khoảng năm lần, ổn định qua năm
mô hình và gần như không bị ảnh hưởng khi đổi tên định danh. Cải tiến CovGate gắn độ phủ vào pipeline chỉ giúp ít.

Thứ ba, đánh giá định vị lỗi bằng LLM nhạy với cách chấm và với bộ dữ liệu: cách phá hòa, cách tính mẫu số và các dòng
lỗi chuẩn sai đều làm thay đổi con số đáng kể.
\IfAtLeast{lan2}{Phân tích thất bại cho thấy thêm rằng độ chính xác cấp dòng phụ thuộc vào việc dòng lỗi được gán ở dòng
điều kiện hay ở câu lệnh mà điều kiện bảo vệ.}
Vì vậy, nhóm đề xuất mọi đánh giá như vậy kèm bảng kiểm tra sáu mục: phá hòa bằng hạng trung bình, mẫu số cố định, dòng
lỗi chuẩn được kiểm chứng, loại bỏ từng thành phần kể cả thành phần ẩn, lặp và kiểm định, báo cáo chi phí.

\section{Hướng phát triển}
\label{sec:huong-phat-trien}

\IfAtLeast{full}{}{\textbf{Kế hoạch đến bài nộp cuối.}
\IfAtLeast{lan2}{Hoàn thiện bài báo tiếng Anh theo chuẩn arXiv, tự phản biện theo năm chiều, chỉnh sửa theo góp ý của
giảng viên sau hai lần báo cáo và chuẩn bị slide bảo vệ.}{Phân tích các ca thất bại và phân nhóm nguyên nhân, chốt
danh sách đóng góp và quan điểm khoa học cùng các mối đe dọa đến tính hợp lệ (báo cáo lần 2), sau đó hoàn thiện bài báo
tiếng Anh theo chuẩn arXiv.}

}Về nghiên cứu, ba hướng mở ra từ kết quả: (i) oracle kiểm chứng dòng lỗi cho phép nhiều vị trí và chống các bản sửa
``bù trừ'' (làm kiểm thử đạt mà không sửa đúng chỗ lỗi), kèm phân xử thủ công các ca tranh chấp; (ii) đánh giá trên lỗi
thật của nhiều dự án hơn và ở định vị cấp hàm trước khi xuống cấp dòng; (iii) một chuẩn đánh giá chung cho định vị lỗi
bằng LLM, trong đó baseline ``một prompt có bằng chứng thực thi'' là mốc bắt buộc.
